\documentclass[11pt]{article}
\usepackage[T1]{fontenc}
\usepackage{lmodern}
\usepackage[margin=1.13in]{geometry}
\usepackage{microtype}
\usepackage{fancyvrb}
\usepackage{amsmath,amssymb,amsthm}
\usepackage[hidelinks]{hyperref}
\hypersetup{pdftitle={Ellipsoids disprove the covering-product conjecture}}
\newtheorem{theorem}{Theorem}
\newtheorem{lemma}[theorem]{Lemma}
\newtheorem{corollary}[theorem]{Corollary}
\newcommand{\B}{\mathbb B}
\newcommand{\R}{\mathbb R}
\newcommand{\Z}{\mathbb Z}
\newcommand{\vol}{\operatorname{vol}}
\newcommand{\rk}{\operatorname{rank}}
\newcommand{\proj}{\operatorname{proj}}
\newcommand{\Pcov}{\mathcal P}
\title{Ellipsoids disprove the covering-product conjecture}
\author{}
\date{27 November 2026}
\begin{document}
\maketitle

\begin{abstract}
Gonz\'alez Merino and Schymura conjectured that the volume of a convex
body times the product of its covering minima with respect to
\(\mathbb Z^n\) is at least \((n+1)/2^n\).
We construct origin-centered ellipsoids for which this product is at most
\(\exp(-\tfrac14 n\log n+O(n))\).
The proof passes from simultaneous covolume bounds for all ranks in
a Haar-random dual lattice to covering-radius bounds for every quotient.
In particular the conjecture fails in every sufficiently large dimension.
A separate exact-arithmetic certificate proves the failure at dimension
\(32768\).
\end{abstract}

\section{The covering-product conjecture}

For a convex body \(K\subset\R^n\) and a full-rank lattice \(\Lambda\),
let \(\mu_i(K,\Lambda)\) be the infimum of \(t>0\) for which
\(tK+\Lambda\) meets every affine subspace of dimension \(n-i\).
The case \(i=n\) is the covering radius of \(K\) with respect to
\(\Lambda\). Define the normalized covering product
\[
 \Pcov(K,\Lambda)=
 \frac{\vol_n(K)}{\det\Lambda}\prod_{i=1}^n\mu_i(K,\Lambda).
\]
Here a lattice determinant is its Euclidean covolume in its own span.
All logarithms are natural. Write \(\B^n\) for the Euclidean unit ball
and \(V_n=\vol_n(\B^n)\).

Gonz\'alez Merino and Schymura proved the lower bound
\(\Pcov(K,\Z^n)\ge1/n!\) and conjectured the stronger inequality
\begin{equation}\label{eq:conj}
 \Pcov(K,\Z^n)\ge\frac{n+1}{2^n}
 \qquad\text{for every convex body }K\subset\R^n.
\end{equation}
This is Conjecture 4.8 of the published paper \cite{GMS},
Conjecture 4.7 of its preprint \cite{GMH}, and Conjecture F of
Codenotti, Santos and Schymura \cite[Sec.~5.2]{CSS}.
The proposed equality example is the standard terminal simplex.
The sharp planar inequality and the bounds for weighted terminal
simplices are consistent with \eqref{eq:conj}. The conjecture
is also recorded by Krivoku\'{c}a \cite{Krivokuca} and in work on \(\kappa\)-covering minima
\cite[Sec.~1]{Xue}, which impose a common-frame condition.
The covering minima in \eqref{eq:conj} are the classical ones.

\begin{theorem}\label{thm:main}
For every integer \(n\ge2\) there exists an origin-centered ellipsoid
\(E\subset\R^n\) of volume \(V_n\) such that for every \(1\le i\le n\),
\begin{equation}\label{eq:profile}
 \mu_i(E,\Z^n)
 \le e^4\sqrt{2e/\pi}\,\sqrt i\ n^{-(n-i)/(2n)}.
\end{equation}
Consequently
\begin{equation}\label{eq:result}
 \Pcov(E,\Z^n)
 \le \bigl(5e^4\sqrt{2e/\pi}\bigr)^n n^{-(n-1)/4}.
\end{equation}
For every integer \(n\ge 2^{44}\), this upper bound is less
than \(2^{-n}\). The conjecture also fails at \(n=32768\).
\end{theorem}

The proof treats the first \(i\) ranks of a dual lattice as one
condition and converts that condition into a bound for the \(i\)-th
covering minimum. Using only a bound on the ambient covering radius
would omit the factor \(n^{-(n-i)/(2n)}\) at rank \(i\).
The product of these factors is \(n^{-(n-1)/4}\).

\section{A simultaneous covolume profile}

A sublattice \(M\subset L\) is primitive if
\(M=L\cap\operatorname{span}_{\R}(M)\). Define
\[
 \delta_k(L)=\inf_{\substack{M\subset L\\\rk M=k}}\det M
 \qquad (1\le k\le n).
\]
Saturating a sublattice in its span decreases its determinant by an
integer index. It therefore suffices to bound primitive sublattices.

\begin{lemma}\label{lem:profile}
For \(n\ge2\), the Haar probability measure on the space of
unimodular lattices gives probability at least
\(1-(e^n-1)^{-1}\) to the simultaneous inequalities
\begin{equation}\label{eq:deltas}
 \delta_k(L)^{1/k}\ge e^{-4}n^{(n-k)/(2n)}
 \qquad(1\le k\le n).
\end{equation}
\end{lemma}

\begin{proof}
For \(1\le k<n\) let \(N_k(L,H)\) count primitive rank-\(k\)
sublattices of determinant at most \(H\).
The one-sublattice case of Kim's mean value theorem
\cite[Thm.~3]{Kim}, previously obtained by Thunder, is
\[
 {\mathbb E}N_k(L,H)=a(n,k)H^n,\qquad
 a(n,k)=\frac1n\binom nk\prod_{j=1}^k
       \frac{V_{n-j+1}\zeta(j)}{V_j\zeta(n-j+1)}.
\]
Here \(\zeta(1)\) is set to 1 in this formula.
For integer \(1\le j\le n\), the inscribed cube and the Gaussian
integral give
\[
 (2/\sqrt j)^j\le V_j\le \bigl(\sqrt{2\pi e/j}\bigr)^j .
\]
For the upper bound, integrate \(e^{-j\|x\|^2/2}\) over
\(\R^j\), using its minimum \(e^{-j/2}\) on the unit ball.
It follows that
\[
 0\le\log V_j+\frac j2\log n<2n .
\]
Indeed the lower bound gives at least
\(j\log2+\tfrac j2\log(n/j)\), and the upper bound gives at most
\(\tfrac n2\log(2\pi e)+\tfrac n{2e}<2n\).
We used \(u\log(1/u)\le1/e\) for \(0<u\le1\).
Thus each logarithmic volume ratio in \(a(n,k)\) is at most
\(-\tfrac12(n-2j+1)\log n+2n\).
The zeta ratios are at most 2 and \(\binom nk\le n^k\).
Using \(\log(2n)\le n\), we obtain
\[
 a(n,k)\le e^{3nk}n^{-k(n-k)/2}.
\]
Choose \(H_k=(e^{-4}n^{(n-k)/(2n)})^k\).
The expected total number of primitive sublattices in all ranks
\(k<n\) below the corresponding thresholds is at most
\(\sum_{k=1}^{n-1}e^{-nk}<(e^n-1)^{-1}\).
Markov's inequality and saturation prove the claim
in ranks \(k<n\). Rank \(n\) follows from \(\det L=1\).
\end{proof}

A simultaneous profile of this order already appears in
Shapira and Weiss \cite[Prop.~3]{ShapiraWeiss}. The explicit calculation
above specifies constants and the probability needed here.
The new use of the profile is to retain its rank dependence through
each projected covering problem.

\section{From the dual profile to all covering minima}

For a lattice \(\Gamma\) of full rank in a Euclidean space \(W\),
write \(\mu(\Gamma)=\mu_{\dim W}(\B^n\cap W,\Gamma)\).
We first give the transfer using the explicit bound of Letwin, Reis
and Rothvoss \cite[Thm.~1.2]{LRR}. They show that a lattice
\(\Gamma\) of dimension \(r\) satisfies
\(\mu(\Gamma)\le c_0\sqrt{r}\), with \(c_0=\sqrt{2e/\pi}\),
if every positive-rank sublattice of \(\Gamma^*\) has determinant
at least 1. This follows by applying their bound to the Euclidean
ball and using \(\mathbb E\|G\|_2\le\sqrt r\) for a standard Gaussian
\(G\). An alternative route using published theorems is given below.

\begin{lemma}\label{lem:transfer}
If \(\Lambda\subset\R^n\) is a full-rank lattice and
\[
 d_i=\min_{1\le k\le i}\delta_k(\Lambda^*)^{1/k},
\]
then \(\mu_i(\B^n,\Lambda)\le c_0\sqrt i/d_i\).
\end{lemma}

\begin{proof}
First suppose \(W\) is an \(i\)-dimensional subspace such that
\(\Gamma=\proj_W\Lambda\) is a lattice in \(W\).
Its dual in \(W\) is \(\Lambda^*\cap W\): for \(y\in W\), the pairings
with \(\proj_W\Lambda\) and with \(\Lambda\) coincide.
Every rank-\(k\) sublattice of \(\Gamma^*\) has determinant at least
\(d_i^k\). The dual of \(d_i\Gamma\) is \(d_i^{-1}\Gamma^*\), so it
satisfies the determinant hypothesis. Apply the covering-radius
bound to \(d_i\Gamma\) and divide by \(d_i\).

For completeness, the same bound handles arbitrary directions in
the definition of \(\mu_i\).
For an arbitrary \(i\)-plane \(W\), let \(G=\overline{\proj_W\Lambda}\).
A closed subgroup of a finite-dimensional Euclidean space is the
orthogonal sum \(G=V\oplus\Gamma\) of a vector space and a discrete
lattice in its orthogonal complement.
Here the linear span of \(G\) is \(W\). Thus \(\Gamma\) is
full rank in \(U=W\cap V^\perp\), with \(r=\dim U\le i\).
The vectors \(y\in U\) with integral pairing against \(\Gamma\) are
exactly \(\Lambda^*\cap U\). This follows from continuity of the
pairing and the equality of pairings with \(\Lambda\) and with
\(\proj_W\Lambda\).
The infimum distance from \(\proj_W x\) to \(\proj_W\Lambda\)
is the distance from \(\proj_U x\) to \(\Gamma\).
If \(r=0\), this distance is zero. Otherwise the preceding Gaussian
argument applies to \(\Gamma\) with the same \(d_i\) and dimension
\(r\), so the resulting bound is at most \(c_0\sqrt i/d_i\).
An affine plane \(x+W^\perp\) meets \(t\B^n+\Lambda\) whenever
that infimum distance is strictly less than \(t\).
Taking the infimum in \(t\) now gives the lemma. This is the
projection principle for covering minima of Kannan and Lov\'asz
\cite{KL}.
\end{proof}

\begin{proof}[Proof of Theorem~\ref{thm:main}]
Choose \(L\) as in Lemma~\ref{lem:profile} and set \(\Lambda=L^*\).
Since \(n^{(n-k)/(2n)}\) is decreasing in \(k\), the quantity \(d_i\)
in Lemma~\ref{lem:transfer} is at least
\(e^{-4}n^{(n-i)/(2n)}\).
Substitution gives \eqref{eq:profile} for \((\B^n,\Lambda)\).

Let \(A\) be a basis matrix of \(\Lambda\) and set \(E=A^{-1}\B^n\).
Invertible linear maps preserve the affine incidences defining
the covering minima. Therefore
\(\mu_i(E,\Z^n)=\mu_i(\B^n,\Lambda)\) for each \(i\).
Also \(|\det A|=1\), so \(\vol_n(E)=V_n\).
This puts the construction in the original lattice of the conjecture.

Multiplying the bounds and using
\(V_n\le(5/\sqrt n)^n\) and \(\prod_{i=1}^n\sqrt i\le n^{n/2}\)
gives \eqref{eq:result}, because
\[
 \sum_{i=1}^n \frac{n-i}{2n}=\frac{n-1}{4}.
\]
For \(n\ge 2^{44}\), use \(n^{1/(4n)}\le e^{1/4}<4/3\),
\(e<3\), and \(c_0<3/2\). The \(n\)-th root of the right-hand side
of \eqref{eq:result} is less than
\[
 5\cdot 3^4\cdot (3/2)\cdot(4/3)\cdot 2^{-11}
   =810/2048<1/2 .
\]
Section~\ref{sec:finite} supplies the separate certificate at
dimension \(32768\).
\end{proof}

\begin{corollary}
For a Haar-random unimodular \(\Lambda\subset\R^n\),
\eqref{eq:result} holds for \(\Pcov(\B^n,\Lambda)\) with probability
at least \(1-(e^n-1)^{-1}\).
\end{corollary}
\begin{proof}
Duality preserves Haar probability on unimodular lattices.
The proof of Theorem~\ref{thm:main} uses exactly the event in
Lemma~\ref{lem:profile} for \(L=\Lambda^*\).
\end{proof}

\section{A route using published covering estimates}

The log-free order in \eqref{eq:result} also follows with an
unspecified absolute constant using published work. A determinant-one
lattice is stable if all of its nonzero sublattices have determinant
at least 1. Regev and Stephens-Davidowitz \cite[Thm.~6.7]{RSD},
together with the dimension-independent bound on isotropic constants
of Klartag and Lehec \cite{Slicing}, give
\(\mu(S)\le C\sqrt r\) for every stable rank-\(r\) lattice \(S\)
with an absolute \(C\).

Suppose now that every positive-rank sublattice of \(D\) has
determinant at least 1. Among sublattices \(M\ne 0\), choose a
primitive \(M\) minimizing \(a=(\det M)^{1/\rk M}\). A minimum
exists since bounded determinants in a fixed rank come from
bounded nonzero vectors of the discrete exterior-power lattice.
Saturation lets us choose a primitive minimizer. Then \(a\ge1\)
and \(a^{-1}M\) is stable. Let \(T=\proj_{(\operatorname{span}M)^\perp}D\).
Every rank-\(k\) sublattice of \(T\) has determinant at least \(a^k\):
its preimage in \(D\) has determinant equal to the product of its
determinant with \(\det M\), and satisfies the defining bound for \(a\).
Thus the same hypothesis holds for \(T\).

The dual \(D^*\) has section \(T^*\) in \((\operatorname{span}M)^\perp\)
and orthogonal projection \(M^*\) onto \(\operatorname{span}M\).
This follows by extending a basis for the primitive section \(M\)
to one for \(D\) and taking the dual basis. Choosing a lattice point
first in the projected lattice and then in its fiber gives
\(\mu(D^*)^2\le\mu(M^*)^2+\mu(T^*)^2\) by Pythagoras.
The dual of a stable lattice is stable \cite[Prop.~2.5]{RSD}.
Hence \(M^*=a^{-1}(a^{-1}M)^*\) satisfies
\(\mu(M^*)\le C\sqrt{\rk M}\). Induction on the rank, with the zero
lattice in dimension zero as the base case, gives
\(\mu(D^*)\le C\sqrt{\rk D}\). This replaces \(c_0\) by an absolute
\(C\) in Lemma~\ref{lem:transfer} and gives the claimed log-free order.

The disproof alone can also be obtained from the original reverse
Minkowski Gaussian theorem and its smoothing-to-covering estimate
\cite[Thm.~1.2 and Lemma~6.1]{RSD}. On \(d_i^{-1}\Gamma^*\) the
Gaussian parameter \(1/[10(\log r+2)]\) gives mass at most \(3/2\),
where \(r=\rk\Gamma\le i\). Scaling back and applying the covering
estimate gives \(\mu(\Gamma)\le20(\log i+2)\sqrt i/d_i\).
The extra factor \((\log n+2)^n\) in the product does not remove
the superexponential decay.

\section{An exact certificate at dimension 32,768}\label{sec:finite}

For fixed \(n\), use the exact \(a(n,k)\) in Lemma~\ref{lem:profile}
and take \(H_k=(2n\,a(n,k))^{-1/n}\) for \(1\le k<n\).
The expected total count is \((n-1)/(2n)\), so some unimodular \(L\)
satisfies the corresponding determinant lower bounds. Put
\(h_k=H_k^{1/k}\) for \(k<n\), \(h_n=1\), and
\(b_i=\min_{k\le i}h_k\). Lemma~\ref{lem:transfer} gives
\[
 \log \Pcov(E,\Z^n)
 \le \log V_n+\tfrac n2\log(2e/\pi)
                +\tfrac12\log(n!)+\sum_{i=1}^n\log(1/b_i).
\]
Use the zeta ratio bound \(2\) in \(a(n,k)\). The Stirling bound
with remainder in \((0,1/(6j))\) yields
\[
 \log V_j=
 \tfrac j2\bigl(\log(2\pi)+1-\log j\bigr)
  -\tfrac12(\log\pi+\log j)-R_j,\qquad 0<R_j<\tfrac1{6j}.
\]
For the logarithms in the resulting expression, scale a rational
argument at least 1 by powers of 2 into \([1,2]\). If
\(v=(x-1)/(x+1)\) in this interval, then \(0\le v\le1/3\). The
first 20 terms of \(2\sum_{k\ge0}v^{2k+1}/(2k+1)\) give a lower
bound for \(\log x\), with remainder at most
\(2/(41\cdot 3^{41}(1-1/9))\).
The code below uses these estimates and exact integers with scale
\(10^{12}\). It uses Machin's formula with two alternating-series
brackets for \(\pi\). The resulting upper bound on the log ratio of
the covering product to the conjectured bound is
\[
 -1313273701890227/10^{12} < -1300 .
\]

\begin{Verbatim}[fontsize=\small]
from fractions import Fraction

S = 10**12
N = 32768


def ceil(a, b):
    return -(-a // b)


def log_unit_interval(a, b):
    assert b <= a <= 2 * b
    z_num = a - b
    z_den = a + b
    pn = z_num
    pd = z_den
    lo = 0
    for k in range(20):
        lo += 2 * S * pn // ((2*k+1)*pd)
        pn *= z_num * z_num
        pd *= z_den * z_den
    # Each rounded term errs by <1/S, and the residual is <1/S.
    assert 2*S*9 < 41*8*3**41
    return lo, lo + 21


L2, U2 = log_unit_interval(2, 1)


def log_bounds(a, b=1):
    p = 0
    while a > 2*b:
        b *= 2
        p += 1
    lo, hi = log_unit_interval(a, b)
    return lo + p*L2, hi + p*U2


# Machin formula with exact alternating-series brackets.
def atan_bounds(a, count):
    lo = sum((-1)**k * Fraction(1, (2*k+1)*a**(2*k+1)) for k in range(count))
    assert count % 2 == 0
    return lo, lo+Fraction(1,(2*count+1)*a**(2*count+1))


a,b = atan_bounds(5, 10)
c,d = atan_bounds(239, 4)
assert Fraction(314159,100000) < 16*a-4*d
assert 16*b-4*c < Fraction(314160,100000)
LP = log_bounds(314159,100000)[0]
UP = log_bounds(314160,100000)[1]

L = [0]
U = [0]
FL = [0]
FU = [0]
VL = [0]
VU = [0]
for j in range(1, N+2):
    l,u = log_bounds(j)
    L.append(l)
    U.append(u)
    FL.append(FL[-1]+l)
    FU.append(FU[-1]+u)
    vind_lo = (j*(L2+LP+S-u)-UP-u)//2 - ceil(S,6*j)
    vind_hi = ceil(j*(U2+UP+S-l)-LP-l,2)
    VL.append(vind_lo)
    VU.append(vind_hi)

prefix_m = -10**100
det_penalty = 0
Vsum = 0
for k in range(1, N):
    Vsum += VU[N-k+1]-VL[k]
    loga_upper = FU[N]-FL[k]-FL[N-k]-L[N]+Vsum+k*U2
    logthreshold_hi = ceil(U2+U[N]+loga_upper, N*k)
    prefix_m = max(prefix_m, logthreshold_hi)
    det_penalty += prefix_m

det_penalty += max(prefix_m,0)
log_product_upper = VU[N]+ceil(N*(U2+S-LP)+FU[N],2)+det_penalty
log_conjecture_lower = L[N+1]-N*U2
print('dimension:', N)
print('exact upper numerator:', log_product_upper-log_conjecture_lower)
print('exact upper denominator:', S)
assert log_product_upper-log_conjecture_lower < -1300*S
\end{Verbatim}

\section{Relation to earlier results}

The corpus manuscripts \cite{VMRC200,VMRC221,VMRC225} examine
covering minima by quotient reduction, exact simplex formulas, and
closure under products and monotone gauge aggregation.
The last two leave the unrestricted covering-product conjecture open.
Lovett and Regev \cite{LovettRegev} combined the random lattice profile
with reverse Minkowski to bound lattice-point counts in sublattices
for a different conjecture. Here the auxiliary problem is to control every quotient dimension
from a single dual covolume profile. Its solution uses the
reverse Minkowski Gaussian bound separately in each dimension.
The quotient estimates multiply to give the superexponential
bound for the original product in the fixed integer lattice.

The disproof concerns the universal inequality \eqref{eq:conj}.
It leaves the sharp lower bound for the classical covering product
and the least counterexample dimension undetermined.
Strict inequality also yields rational ellipsoid and rational
centrally symmetric polytope counterexamples. For example, choose
a rational polytope \(P\) with \((1-\varepsilon)E\subseteq P
\subseteq(1+\varepsilon)E\). Such symmetric polytopes follow by
taking a sufficiently fine finite net on the ellipsoid, approximating
the points by rational pairs \(v,-v\), and taking the convex hull.
Monotonicity of volume and reverse monotonicity of every covering
minimum give
\(\Pcov(P,\Z^n)\le((1+\varepsilon)/(1-\varepsilon))^n\Pcov(E,\Z^n)\).
Taking \(\varepsilon\) small preserves the strict inequality.
The same argument applies to rational positive-definite quadratic
forms close to the form defining \(E\).

The literature comparison includes the original conjecture, the
subsequent simplex bounds, the simultaneous covolume profile, the
reverse Minkowski theorem, and the recent \(\kappa\)-covering
product result. No earlier ellipsoid disproof was identified in
these sources or the targeted literature search. This does not
establish priority against material outside the search.

\small
\begin{thebibliography}{99}
\setlength{\itemsep}{3pt}

\bibitem{CSS}
G.~Codenotti, F.~Santos and M.~Schymura,
The covering radius and a discrete surface area for non-hollow simplices,
\emph{Discrete Comput. Geom.} \textbf{67} (2022), 65--111.
\href{https://doi.org/10.1007/s00454-021-00330-3}{DOI: 10.1007/s00454-021-00330-3}.

\bibitem{GMH}
B.~Gonz\'alez Merino and M.~Henze,
On densities of lattice arrangements intersecting every \(i\)-dimensional
affine subspace, \href{https://arxiv.org/abs/1605.00443}{arXiv:1605.00443}, 2016.

\bibitem{GMS}
B.~Gonz\'alez Merino and M.~Schymura,
On densities of lattice arrangements intersecting every \(i\)-dimensional
affine subspace, \emph{Discrete Comput. Geom.} \textbf{58} (2017),
663--685.
\href{https://doi.org/10.1007/s00454-017-9911-x}{DOI: 10.1007/s00454-017-9911-x}.

\bibitem{Krivokuca}
K.~Krivoku\'{c}a, Upper bounds on covering minima of convex bodies,
\href{https://arxiv.org/abs/2601.15173}{arXiv:2601.15173}, 2026.

\bibitem{KL}
R.~Kannan and L.~Lov\'asz, Covering minima and lattice-point-free convex
bodies, \emph{Ann. of Math.} \textbf{128} (1988), 577--602.
\href{https://doi.org/10.2307/1971436}{DOI: 10.2307/1971436}.

\bibitem{Kim}
S.~Kim, Mean value formulas on sublattices and flags of the random lattice,
\emph{J. Number Theory}, 2022.
\href{https://doi.org/10.1016/j.jnt.2022.03.013}{DOI: 10.1016/j.jnt.2022.03.013}.
\href{https://arxiv.org/abs/2005.10874}{arXiv:2005.10874v3}.

\bibitem{LRR}
B.~Letwin, V.~Reis and T.~Rothvoss, Strongly log-concave measures and
lattices, \href{https://arxiv.org/abs/2610.29238}{arXiv:2610.29238}, 2026.

\bibitem{LovettRegev}
S.~Lovett and O.~Regev, A counterexample to a strong variant of the
polynomial Freiman--Ruzsa conjecture in Euclidean space,
\emph{Discrete Analysis} 2017:8.
\href{https://doi.org/10.19086/da.1640}{DOI: 10.19086/da.1640}.

\bibitem{RSD}
O.~Regev and N.~Stephens-Davidowitz, A reverse Minkowski theorem,
\emph{Ann. of Math.} \textbf{199} (2024), 1--49.
\href{https://doi.org/10.4007/annals.2024.199.1.1}{DOI: 10.4007/annals.2024.199.1.1}.

\bibitem{ShapiraWeiss}
U.~Shapira and B.~Weiss, A volume estimate for the set of stable
lattices, \emph{C. R. Math. Acad. Sci. Paris} \textbf{352} (2014),
875--879.
\href{https://doi.org/10.1016/j.crma.2014.08.019}{DOI: 10.1016/j.crma.2014.08.019}.

\bibitem{Slicing}
B.~Klartag and J.~Lehec, Affirmative resolution of Bourgain's slicing
problem using Guan's bound, \emph{Geom. Funct. Anal.}
\textbf{35} (2025), 1147--1168.
\href{https://doi.org/10.1007/s00039-025-00718-w}{DOI: 10.1007/s00039-025-00718-w}.

\bibitem{VMRC200}
Verifiable Mathematical Reasoning Corpus, manuscript [200],
\emph{Weighted fractional-part flows and sharp lattice covering bounds}, snapshot
\href{https://github.com/QEEDDD/Verifiable-Mathematical-Reasoning-Corpus/tree/e73bae9fd0c319129b19aacc5c317f7abf2865ef}{e73bae9f},
2026.

\bibitem{VMRC221}
Verifiable Mathematical Reasoning Corpus, manuscript [221],
\emph{Further closures in lattice covering},
snapshot e73bae9f, 2026.

\bibitem{VMRC225}
Verifiable Mathematical Reasoning Corpus, manuscript [225],
\emph{Arithmetic gaps, exact minima, and covering products},
snapshot e73bae9f, 2026.

\bibitem{Xue}
F.~Xue, On \(\kappa\)-covering minima,
\href{https://arxiv.org/abs/2610.39034}{arXiv:2610.39034}, 2026.

\end{thebibliography}
\end{document}
